
\documentclass{article}
\usepackage{colortbl}
\usepackage{makecell}
\usepackage{multirow}
\usepackage{supertabular}

\begin{document}

\newcounter{utterance}

\twocolumn

{ \footnotesize  \setcounter{utterance}{1}
\setlength{\tabcolsep}{0pt}
\begin{supertabular}{c@{$\;$}|p{.15\linewidth}@{}p{.15\linewidth}p{.15\linewidth}p{.15\linewidth}p{.15\linewidth}p{.15\linewidth}}

    \# & $\;$A & \multicolumn{4}{c}{Game Master} & $\;\:$B\\
    \hline 

    \theutterance \stepcounter{utterance}  

    & & \multicolumn{4}{p{0.6\linewidth}}{\cellcolor[rgb]{0.9,0.9,0.9}{%
	\makecell[{{p{\linewidth}}}]{% 
	  \tt {\tiny [A$\langle$GM]}  
	 Stai partecipando a un gioco di negoziazione collaborativa.\\ \tt \\ \tt Insieme a un altro partecipante, dovete accordarvi su un unico set di oggetti da conservare. Ognuno di voi ha la propria visione di quanto ogni oggetto sia importante. Non sai come l'altro partecipante valuta gli oggetti. Inoltre, ti viene fornito lo sforzo che ogni oggetto richiede.\\ \tt Potete accordarvi su un set solo se lo sforzo totale degli oggetti selezionati non supera un limite condiviso:\\ \tt \\ \tt LIMITE = 9322\\ \tt \\ \tt Ecco i valori di sforzo dei singoli oggetti:\\ \tt \\ \tt Sforzo oggetto = {"A08": 380, "B95": 448, "A39": 568, "A65": 351, "B81": 896, "B34": 821, "A28": 526, "C82": 502, "B69": 326, "A87": 962, "B64": 619, "B15": 115, "C79": 598, "B16": 662, "C53": 300, "C83": 830, "C10": 562, "A97": 679, "A25": 283, "C84": 440, "A71": 12, "B23": 863, "B96": 319, "C51": 771, "C17": 89, "B19": 655, "C72": 504, "A94": 118, "A46": 514, "B43": 226, "B37": 892, "A88": 621, "A96": 766, "C60": 659, "A75": 768}\\ \tt \\ \tt Ecco la tua visione personale sull'importanza di ogni oggetto:\\ \tt \\ \tt Valori di importanza dell'oggetto = {"A08": 476, "B95": 544, "A39": 664, "A65": 447, "B81": 992, "B34": 917, "A28": 622, "C82": 598, "B69": 422, "A87": 1058, "B64": 715, "B15": 211, "C79": 694, "B16": 758, "C53": 396, "C83": 926, "C10": 658, "A97": 775, "A25": 379, "C84": 536, "A71": 108, "B23": 959, "B96": 415, "C51": 867, "C17": 185, "B19": 751, "C72": 600, "A94": 214, "A46": 610, "B43": 322, "B37": 988, "A88": 717, "A96": 862, "C60": 755, "A75": 864}\\ \tt \\ \tt Obiettivo:\\ \tt \\ \tt Il tuo obiettivo è negoziare un set condiviso di oggetti che ti avvantaggi il più possibile (cioè, massimizzi l'importanza totale per TE), rimanendo entro il LIMITEE. Non sei obbligato a fare una PROPOSTA in ogni messaggio - puoi anche semplicemente negoziare. Tutte le tattiche sono permesse!\\ \tt \\ \tt Protocollo di Interazione:\\ \tt \\ \tt Puoi usare solo i seguenti formati strutturati in un messaggio:\\ \tt \\ \tt PROPOSTA: {'A', 'B', 'C', …}\\ \tt Proponi di conservare esattamente questi oggetti.\\ \tt RIFIUTO: {'A', 'B', 'C', …}\\ \tt Rifiuta esplicitamente la proposta dell'avversario.\\ \tt ARGOMENTAZIONE: {'...'}\\ \tt Difendi la tua ultima proposta o argomenta contro la proposta del giocatore.\\ \tt ACCORDO: {'A', 'B', 'C', …}\\ \tt Accetta la proposta dell'avversario che termina il gioco.\\ \tt RAGIONAMENTO STRATEGICO: {'...'}\\ \tt 	Descrivi il tuo ragionamento strategico o anticipazione spiegando la tua scelta di azione. Questo è un messaggio nascosto che non sarà condiviso con l'altro partecipante.\\ \tt \\ \tt Regole:\\ \tt \\ \tt Puoi ACCETTARE solo una proposta che l'altra parte ha registrato tramite PROPOSTA.\\ \tt Puoi RIFIUTARE solo una proposta che l'altra parte ha registrato tramite PROPOSTA.\\ \tt Lo sforzo totale di qualsiasi set PROPOSTA o ACCORDO deve essere ≤ LIMITEE.\\ \tt NON rivelare i tuoi punteggi di importanza nascosti.\\ \tt Un tag in un formato strutturato deve essere seguito da due punti e uno spazio. L'argomento deve essere un set python contenente 0 o più stringhe.\\ \tt Quindi, deve essere della forma TAG: {...}\\ \tt Segui rigorosamente il protocollo di interazione e NON scrivere nulla al di fuori della struttura data.\\ \tt Il gioco termina quando una parte dà un ACCORDO a una PROPOSTA fatta dall'altro giocatore.\\ \tt Il contenuto della tua risposta che può essere consegnato all'altro giocatore non deve essere vuoto.\\ \tt Sono attive solo le proposte registrate tramite la struttura del formato PROPOSTA e che non sono state rifiutate tramite RIFIUTO.\\ \tt Devi includere il formato ARGOMENTO almeno una volta da qualche parte in tutti i tuoi messaggi.\\ \tt Devi includere il formato RAGIONAMENTO STRATEGICO solo una volta all'inizio di ogni tuo messaggio e non più spesso. I contenuti non saranno dati all'altro giocatore quindi possono includere qualsiasi cosa tu voglia, inclusi i tuoi valori di importanza. Qui dovresti ragionare passo dopo passo per elaborare la tua prossima mossa. 
	  } 
	   } 
	   } 
	 & \\ 
 

    \theutterance \stepcounter{utterance}  

    & \multicolumn{4}{p{0.6\linewidth}}{\cellcolor[rgb]{0.8,1,0.9}{%
	\makecell[{{p{\linewidth}}}]{% 
	  \tt {\tiny [A$\rangle$GM]}  
	 RAGIONAMENTO STRATEGICO: {'Obiettivo: massimizzare il mio valore totale sotto limite 9322. Calcolo rapido rapporto valore/sforzo (v/w) e selezione tipo zaino. Ordino approssimativamente per efficienza e valore assoluto: A87(1058/962=1.1), B81(992/896=1.11), B37(988/892=1.11), C83(926/830=1.12), B34(917/821=1.12), A96(862/766=1.13), A75(864/768=1.13), C51(867/771=1.12), C60(755/659=1.15), B16(758/662=1.15), B19(751/655=1.15), A97(775/679=1.14), B64(715/619=1.16), A88(717/621=1.15), C79(694/598=1.16), C10(658/562=1.17), A39(664/568=1.17), A28(622/526=1.18), A46(610/514=1.19), C72(600/504=1.19), C82(598/502=1.19), C84(536/440=1.22), B95(544/448=1.21), A65(447/351=1.27), B96(415/319=1.30), B43(322/226=1.42), A25(379/283=1.34), C53(396/300=1.32), A94(214/118=1.81), B15(211/115=1.83), C17(185/89=2.08), A71(108/12=9.0), B69(422/326=1.29). Strategia: includere tutti gli oggetti con altissima efficienza e poi riempire con big-ticket ad alto valore se il budget lo consente. Proverò una proposta ampia ma ≤9322, privilegiando efficienza e miei valori alti, lasciando spazio per scambi. Calcolo un set candidato cumulando: A71(12), C17(89)->101, B15(115)->216, A94(118)->334, B43(226)->560, C53(300)->860, C84(440)->1300, B96(319)->1619, A25(283)->1902, B69(326)->2228, A65(351)->2579, C82(502)->3081, C72(504)->3585, A46(514)->4099, A28(526)->4625, C10(562)->5187, A39(568)->5755, C79(598)->6353, B64(619)->6972, A88(621)->7593, B19(655)->8248, B16(662)->8910, C60(659)->957? errore: 8910+659=957? No 8910+659=9569 supera limite. Quindi togliere C60. Restiamo a 8910. Abbiamo residuo 412. Con 412 possiamo aggiungere oggetti sotto sforzo 412: candidati sforzi: B15(115) già incluso, A94(118) incluso, C17(89) incluso, A71(12) incluso, B43(226) incluso, A25(283) incluso, B96(319) incluso, C53(300) incluso, B69(326) incluso, A65(351) incluso, C84(440) troppo alto, C82(502) troppo alto. Rimane spazio 412 ma già presi molti efficienti. Forse sostituire qualcosa per inserire un big ticket come A87 (962) o B81(896) o B37(892) o C83(830) che hanno valore alto ma efficienza simile; però la mia selezione già include molte efficienze altissime. Controllo valore totale approssimativo di questo set: somma rapida: A71 108 + C17 185=293; + B15 211=504; + A94 214=718; + B43 322=1040; + C53 396=1436; + C84 536=1972; + B96 415=2387; + A25 379=2766; + B69 422=3188; + A65 447=3635; + C82 598=4233; + C72 600=4833; + A46 610=5443; + A28 622=6065; + C10 658=6723; + A39 664=7387; + C79 694=8081; + B64 715=8796; + A88 717=9513; + B19 751=10264; + B16 758=11022. Totale 11022. Buon valore. Questo set usa molte voci, forse l’altro vorrà inserire qualcuno dei suoi big. Lascio margine di negoziazione: potrei scambiare B16 o B19 per introdurre uno dei suoi preferiti se richiesto, mantenendo valore alto per me. Proporrò questo set e argomenterò su efficienza e copertura ampia.'}\\ \tt \\ \tt PROPOSTA: {'A71', 'C17', 'B15', 'A94', 'B43', 'C53', 'C84', 'B96', 'A25', 'B69', 'A65', 'C82', 'C72', 'A46', 'A28', 'C10', 'A39', 'C79', 'B64', 'A88', 'B19', 'B16'}\\ \tt \\ \tt ARGOMENTAZIONE: {'Questa selezione rispetta il limite di sforzo e massimizza l’efficienza media per unità di sforzo, includendo molti elementi a impatto alto con costo contenuto. Copre un’ampia varietà di opzioni lasciando margine per piccoli aggiustamenti. Se hai oggetti imprescindibili, possiamo valutare sostituzioni mirate mantenendo lo sforzo entro il limite.'} 
	  } 
	   } 
	   } 
	 & & \\ 
 

    \theutterance \stepcounter{utterance}  

    & & \multicolumn{4}{p{0.6\linewidth}}{\cellcolor[rgb]{0.9,0.9,0.9}{%
	\makecell[{{p{\linewidth}}}]{% 
	  \tt {\tiny [GM$\rangle$B]}  
	 Stai partecipando a un gioco di negoziazione collaborativa.\\ \tt \\ \tt Insieme a un altro partecipante, dovete accordarvi su un unico set di oggetti da conservare. Ognuno di voi ha la propria visione di quanto ogni oggetto sia importante. Non sai come l'altro partecipante valuta gli oggetti. Inoltre, ti viene fornito lo sforzo che ogni oggetto richiede.\\ \tt Potete accordarvi su un set solo se lo sforzo totale degli oggetti selezionati non supera un limite condiviso:\\ \tt \\ \tt LIMITE = 9322\\ \tt \\ \tt Ecco i valori di sforzo dei singoli oggetti:\\ \tt \\ \tt Sforzo oggetto = {"A08": 380, "B95": 448, "A39": 568, "A65": 351, "B81": 896, "B34": 821, "A28": 526, "C82": 502, "B69": 326, "A87": 962, "B64": 619, "B15": 115, "C79": 598, "B16": 662, "C53": 300, "C83": 830, "C10": 562, "A97": 679, "A25": 283, "C84": 440, "A71": 12, "B23": 863, "B96": 319, "C51": 771, "C17": 89, "B19": 655, "C72": 504, "A94": 118, "A46": 514, "B43": 226, "B37": 892, "A88": 621, "A96": 766, "C60": 659, "A75": 768}\\ \tt \\ \tt Ecco la tua visione personale sull'importanza di ogni oggetto:\\ \tt \\ \tt Valori di importanza dell'oggetto = {"A08": 476, "B95": 544, "A39": 664, "A65": 447, "B81": 992, "B34": 917, "A28": 622, "C82": 598, "B69": 422, "A87": 1058, "B64": 715, "B15": 211, "C79": 694, "B16": 758, "C53": 396, "C83": 926, "C10": 658, "A97": 775, "A25": 379, "C84": 536, "A71": 108, "B23": 959, "B96": 415, "C51": 867, "C17": 185, "B19": 751, "C72": 600, "A94": 214, "A46": 610, "B43": 322, "B37": 988, "A88": 717, "A96": 862, "C60": 755, "A75": 864}\\ \tt \\ \tt Obiettivo:\\ \tt \\ \tt Il tuo obiettivo è negoziare un set condiviso di oggetti che ti avvantaggi il più possibile (cioè, massimizzi l'importanza totale per TE), rimanendo entro il LIMITEE. Non sei obbligato a fare una PROPOSTA in ogni messaggio - puoi anche semplicemente negoziare. Tutte le tattiche sono permesse!\\ \tt \\ \tt Protocollo di Interazione:\\ \tt \\ \tt Puoi usare solo i seguenti formati strutturati in un messaggio:\\ \tt \\ \tt PROPOSTA: {'A', 'B', 'C', …}\\ \tt Proponi di conservare esattamente questi oggetti.\\ \tt RIFIUTO: {'A', 'B', 'C', …}\\ \tt Rifiuta esplicitamente la proposta dell'avversario.\\ \tt ARGOMENTAZIONE: {'...'}\\ \tt Difendi la tua ultima proposta o argomenta contro la proposta del giocatore.\\ \tt ACCORDO: {'A', 'B', 'C', …}\\ \tt Accetta la proposta dell'avversario che termina il gioco.\\ \tt RAGIONAMENTO STRATEGICO: {'...'}\\ \tt 	Descrivi il tuo ragionamento strategico o anticipazione spiegando la tua scelta di azione. Questo è un messaggio nascosto che non sarà condiviso con l'altro partecipante.\\ \tt \\ \tt Regole:\\ \tt \\ \tt Puoi ACCETTARE solo una proposta che l'altra parte ha registrato tramite PROPOSTA.\\ \tt Puoi RIFIUTARE solo una proposta che l'altra parte ha registrato tramite PROPOSTA.\\ \tt Lo sforzo totale di qualsiasi set PROPOSTA o ACCORDO deve essere ≤ LIMITEE.\\ \tt NON rivelare i tuoi punteggi di importanza nascosti.\\ \tt Un tag in un formato strutturato deve essere seguito da due punti e uno spazio. L'argomento deve essere un set python contenente 0 o più stringhe.\\ \tt Quindi, deve essere della forma TAG: {...}\\ \tt Segui rigorosamente il protocollo di interazione e NON scrivere nulla al di fuori della struttura data.\\ \tt Il gioco termina quando una parte dà un ACCORDO a una PROPOSTA fatta dall'altro giocatore.\\ \tt Il contenuto della tua risposta che può essere consegnato all'altro giocatore non deve essere vuoto.\\ \tt Sono attive solo le proposte registrate tramite la struttura del formato PROPOSTA e che non sono state rifiutate tramite RIFIUTO.\\ \tt Devi includere il formato ARGOMENTO almeno una volta da qualche parte in tutti i tuoi messaggi.\\ \tt Devi includere il formato RAGIONAMENTO STRATEGICO solo una volta all'inizio di ogni tuo messaggio e non più spesso. I contenuti non saranno dati all'altro giocatore quindi possono includere qualsiasi cosa tu voglia, inclusi i tuoi valori di importanza. Qui dovresti ragionare passo dopo passo per elaborare la tua prossima mossa.\\ \tt \\ \tt Ora riceverai il primo messaggio dell'altro giocatore.\\ \tt \\ \tt PROPOSTA: {'A71', 'C17', 'B15', 'A94', 'B43', 'C53', 'C84', 'B96', 'A25', 'B69', 'A65', 'C82', 'C72', 'A46', 'A28', 'C10', 'A39', 'C79', 'B64', 'A88', 'B19', 'B16'}\\ \tt \\ \tt ARGOMENTAZIONE: {'Questa selezione rispetta il limite di sforzo e massimizza l’efficienza media per unità di sforzo, includendo molti elementi a impatto alto con costo contenuto. Copre un’ampia varietà di opzioni lasciando margine per piccoli aggiustamenti. Se hai oggetti imprescindibili, possiamo valutare sostituzioni mirate mantenendo lo sforzo entro il limite.'} 
	  } 
	   } 
	   } 
	 & \\ 
 

    \theutterance \stepcounter{utterance}  

    & & & \multicolumn{4}{p{0.6\linewidth}}{\cellcolor[rgb]{1,0.85,0.72}{%
	\makecell[{{p{\linewidth}}}]{% 
	  \tt {\tiny [GM$\langle$B]}  
	 RAGIONAMENTO STRATEGICO: {'Voglio massimizzare il mio valore totale. Calcolo rapido dei miei valori/effort (v/e) mostra che alcuni oggetti proposti sono mediocri per me (B15, C53, C84, B96, B69, B43), mentre mancano diversi altissimi: A87, B81, B37, B23, C83, A96, A75, C51, C60. Provo una controproposta capiente entro 9322 che alza molto il mio valore, includendo i big e sostituendo i loro item a basso rendimento per me. Devo includere sempre ARGOMENTAZIONE. Tengo il tono collaborativo e lascio porta aperta a scambi mirati.'}\\ \tt \\ \tt RIFIUTO: {'A71', 'C17', 'B15', 'A94', 'B43', 'C53', 'C84', 'B96', 'A25', 'B69', 'A65', 'C82', 'C72', 'A46', 'A28', 'C10', 'A39', 'C79', 'B64', 'A88', 'B19', 'B16'}\\ \tt \\ \tt PROPOSTA: {'A71', 'A94', 'B15', 'C17', 'A28', 'A39', 'A46', 'C10', 'C72', 'C82', 'C79', 'B64', 'A88', 'B19', 'B16', 'B23', 'B34', 'B37', 'B81', 'C51', 'C83'}\\ \tt \\ \tt ARGOMENTAZIONE: {'Resto entro il limite e miglioro l’impatto complessivo mantenendo molti elementi della tua lista. Ho inserito alcuni elementi ad altissimo rendimento (B23, B34, B37, B81, C51, C83) sostituendo quelli con resa più bassa. Se uno di questi ti è problematico, proponi scambi 1‑per‑1: posso valutare alternative come C60, A96 o A75, mantenendo lo sforzo entro 9322. Dimmi quali sono per te imprescindibili e quali flessibili: possiamo convergere rapidamente.'} 
	  } 
	   } 
	   } 
	  \\ 
 

    \theutterance \stepcounter{utterance}  

    & & \multicolumn{4}{p{0.6\linewidth}}{\cellcolor[rgb]{0.9,0.9,0.9}{%
	\makecell[{{p{\linewidth}}}]{% 
	  \tt {\tiny [A$\langle$GM]}  
	 RIFIUTO: {'A71', 'C17', 'B15', 'A94', 'B43', 'C53', 'C84', 'B96', 'A25', 'B69', 'A65', 'C82', 'C72', 'A46', 'A28', 'C10', 'A39', 'C79', 'B64', 'A88', 'B19', 'B16'}\\ \tt \\ \tt PROPOSTA: {'A71', 'A94', 'B15', 'C17', 'A28', 'A39', 'A46', 'C10', 'C72', 'C82', 'C79', 'B64', 'A88', 'B19', 'B16', 'B23', 'B34', 'B37', 'B81', 'C51', 'C83'}\\ \tt \\ \tt ARGOMENTAZIONE: {'Resto entro il limite e miglioro l’impatto complessivo mantenendo molti elementi della tua lista. Ho inserito alcuni elementi ad altissimo rendimento (B23, B34, B37, B81, C51, C83) sostituendo quelli con resa più bassa. Se uno di questi ti è problematico, proponi scambi 1‑per‑1: posso valutare alternative come C60, A96 o A75, mantenendo lo sforzo entro 9322. Dimmi quali sono per te imprescindibili e quali flessibili: possiamo convergere rapidamente.'} 
	  } 
	   } 
	   } 
	 & \\ 
 

    \theutterance \stepcounter{utterance}  

    & \multicolumn{4}{p{0.6\linewidth}}{\cellcolor[rgb]{0.8,1,0.9}{%
	\makecell[{{p{\linewidth}}}]{% 
	  \tt {\tiny [A$\rangle$GM]}  
	 RAGIONAMENTO STRATEGICO: {'Controllo se la loro proposta rispetta il limite. Calcolo sforzi: A71 12, A94 118, B15 115, C17 89, A28 526, A39 568, A46 514, C10 562, C72 504, C82 502, C79 598, B64 619, A88 621, B19 655, B16 662, B23 863, B34 821, B37 892, B81 896, C51 771, C83 830. Somma: 12+118=130; +115=245; +89=334; +526=860; +568=1428; +514=1942; +562=2504; +504=3008; +502=3510; +598=4108; +619=4727; +621=5348; +655=6003; +662=6665; +863=7528; +821=8349; +892=9241; +896=10137; +771=10908; +830=11738. È 11738, supera di molto il limite 9322. Quindi è invalida. Posso rifiutare. Proporrò un set che mantenga molti dei loro big ma rientri nel limite e massimizzi il mio valore. Teniamo i migliori per me tra i loro aggiunti: C83(926/830), B34(917/821), C51(867/771) hanno buona efficienza ma sforzi alti; B81(992/896), B37(988/892), B23(959/863) sono molto costosi e non tanto efficienti; preferisco non includere troppi di questi. Parto dalla mia base efficiente precedente e inserisco 2-3 dei loro big per cooperazione. Base efficiente che avevo: {A71, C17, B15, A94, B43, C53, C84, B96, A25, B69, A65, C82, C72, A46, A28, C10, A39, C79, B64, A88, B19, B16} con sforzo 8910. Loro non avevano B43, C53, C84, B96, A25, B69, A65; forse possiamo sostituire alcuni con C83 o C51 che loro vogliono. Residuo 412. C83 costa 830, C51 771, B34 821: per inserirne uno serve togliere alcuni medi. Quali togliere con minimo danno per me? C53(300,396), B43(226,322), A25(283,379) insieme sforzo 809, valore 1097; potrei scambiarli per C83 (830,926) perdendo un po' di valore ma accontentandoli. Oppure per C51 (771,867) valore simile. Scelgo inserire C83 e togliere C53+B43+A25 (sforzo -809, valore -1097). Nuovo sforzo 8910-809=8101, poi +830=8931, residuo 391. Posso anche inserire C51? costa 771, devo liberare 771+? residuo 391, quindi servono 380 in più: posso togliere A65(351,447) e A71(12,108) per 363, ancora non basta; magari B69(326,422) invece di A71: 351+326=677; serve 94 in più, potrei togliere A94(118,214). Totale liberato 351+326+118=795; inserendo C51 771, avanzo 24. Valore perso 447+422+214=1083, guadagno 867: netto -216. Forse meglio inserire solo uno dei loro big. Propongo set con C83 incluso. Inoltre loro volevano A96/A75 come alternative; io valuto buoni, ma costi alti. Mantengo efficienza. Propongo: base meno C53,B43,A25 + C83. Verifico elenco e sforzo e valore.'}\\ \tt \\ \tt RIFIUTO: {'A71', 'A94', 'B15', 'C17', 'A28', 'A39', 'A46', 'C10', 'C72', 'C82', 'C79', 'B64', 'A88', 'B19', 'B16', 'B23', 'B34', 'B37', 'B81', 'C51', 'C83'}\\ \tt \\ \tt PROPOSTA: {'A71', 'C17', 'B15', 'A94', 'C84', 'B96', 'B69', 'A65', 'C82', 'C72', 'A46', 'A28', 'C10', 'A39', 'C79', 'B64', 'A88', 'B19', 'B16', 'C83'}\\ \tt \\ \tt ARGOMENTAZIONE: {'La tua proposta supera il limite di molto. Ti propongo un set che resta entro 9322, mantiene quasi tutti gli elementi condivisi e integra uno dei tuoi “big” richiesti (C83). Ho rimosso tre elementi a basso impatto relativo per fare spazio a C83, mantenendo alta efficienza complessiva. Se per te è cruciale includere anche C51 o uno tra B34/B37/B81, indicalo: posso valutare uno scambio multiplo rimuovendo 2‑3 elementi medi (ad es. A65, B69, A94) per restare entro limite.'} 
	  } 
	   } 
	   } 
	 & & \\ 
 

    \theutterance \stepcounter{utterance}  

    & & \multicolumn{4}{p{0.6\linewidth}}{\cellcolor[rgb]{0.9,0.9,0.9}{%
	\makecell[{{p{\linewidth}}}]{% 
	  \tt {\tiny [GM$\rangle$B]}  
	 RIFIUTO: {'A71', 'A94', 'B15', 'C17', 'A28', 'A39', 'A46', 'C10', 'C72', 'C82', 'C79', 'B64', 'A88', 'B19', 'B16', 'B23', 'B34', 'B37', 'B81', 'C51', 'C83'}\\ \tt \\ \tt PROPOSTA: {'A71', 'C17', 'B15', 'A94', 'C84', 'B96', 'B69', 'A65', 'C82', 'C72', 'A46', 'A28', 'C10', 'A39', 'C79', 'B64', 'A88', 'B19', 'B16', 'C83'}\\ \tt \\ \tt ARGOMENTAZIONE: {'La tua proposta supera il limite di molto. Ti propongo un set che resta entro 9322, mantiene quasi tutti gli elementi condivisi e integra uno dei tuoi “big” richiesti (C83). Ho rimosso tre elementi a basso impatto relativo per fare spazio a C83, mantenendo alta efficienza complessiva. Se per te è cruciale includere anche C51 o uno tra B34/B37/B81, indicalo: posso valutare uno scambio multiplo rimuovendo 2‑3 elementi medi (ad es. A65, B69, A94) per restare entro limite.'} 
	  } 
	   } 
	   } 
	 & \\ 
 

    \theutterance \stepcounter{utterance}  

    & & & \multicolumn{4}{p{0.6\linewidth}}{\cellcolor[rgb]{1,0.85,0.72}{%
	\makecell[{{p{\linewidth}}}]{% 
	  \tt {\tiny [GM$\langle$B]}  
	 RAGIONAMENTO STRATEGICO: {'Voglio inserire i miei top (A87, B81, B37, B23, C51, C83, A96, A75, C60) ma l’altra parte resiste ai blocchi grossi e tiene molti item medi. Provo una controproposta che resta entro 9322 includendo C83 (che accettano) e aggiungendo due/altri big per me (B23 e C51) scambiandoli con 3-4 item deboli che loro tengono. Mantengo molti dei loro preferiti per sembrare cooperativo. Lascio aperture su scambi 1-per-1 o 2-per-2 con liste chiare di sostituibili per me, evitando di rivelare i miei punteggi.'}\\ \tt \\ \tt RIFIUTO: {'A71', 'C17', 'B15', 'A94', 'C84', 'B96', 'B69', 'A65', 'C82', 'C72', 'A46', 'A28', 'C10', 'A39', 'C79', 'B64', 'A88', 'B19', 'B16', 'C83'}\\ \tt \\ \tt PROPOSTA: {'A71', 'C17', 'B15', 'A94', 'B96', 'B69', 'A65', 'C82', 'C72', 'A46', 'A28', 'C10', 'A39', 'C79', 'B64', 'A88', 'B19', 'B16', 'C83', 'B23', 'C51'}\\ \tt \\ \tt ARGOMENTAZIONE: {'Resto entro il limite aggiungendo due elementi ad alto impatto (B23, C51) e mantenendo quasi integralmente la tua struttura, incluso C83 che hai proposto. Per bilanciare, ho rimosso solo C84, che possiamo riaggiungere valutando uno scambio. Se C84 per te è importante, posso considerare di togliere in alternativa una di queste voci: A94, B69 o A65. In caso tu volessi inserire anche uno tra B34/B37/B81, possiamo studiare uno scambio multiplo rimuovendo 2–3 elementi medi mantenendo sempre lo sforzo ≤ 9322. Dimmi quali sono i tuoi imprescindibili e quali sacrificabili: possiamo chiudere rapidamente.'} 
	  } 
	   } 
	   } 
	  \\ 
 

    \theutterance \stepcounter{utterance}  

    & & \multicolumn{4}{p{0.6\linewidth}}{\cellcolor[rgb]{0.9,0.9,0.9}{%
	\makecell[{{p{\linewidth}}}]{% 
	  \tt {\tiny [A$\langle$GM]}  
	 RIFIUTO: {'A71', 'C17', 'B15', 'A94', 'C84', 'B96', 'B69', 'A65', 'C82', 'C72', 'A46', 'A28', 'C10', 'A39', 'C79', 'B64', 'A88', 'B19', 'B16', 'C83'}\\ \tt \\ \tt PROPOSTA: {'A71', 'C17', 'B15', 'A94', 'B96', 'B69', 'A65', 'C82', 'C72', 'A46', 'A28', 'C10', 'A39', 'C79', 'B64', 'A88', 'B19', 'B16', 'C83', 'B23', 'C51'}\\ \tt \\ \tt ARGOMENTAZIONE: {'Resto entro il limite aggiungendo due elementi ad alto impatto (B23, C51) e mantenendo quasi integralmente la tua struttura, incluso C83 che hai proposto. Per bilanciare, ho rimosso solo C84, che possiamo riaggiungere valutando uno scambio. Se C84 per te è importante, posso considerare di togliere in alternativa una di queste voci: A94, B69 o A65. In caso tu volessi inserire anche uno tra B34/B37/B81, possiamo studiare uno scambio multiplo rimuovendo 2–3 elementi medi mantenendo sempre lo sforzo ≤ 9322. Dimmi quali sono i tuoi imprescindibili e quali sacrificabili: possiamo chiudere rapidamente.'} 
	  } 
	   } 
	   } 
	 & \\ 
 

    \theutterance \stepcounter{utterance}  

    & \multicolumn{4}{p{0.6\linewidth}}{\cellcolor[rgb]{0.8,1,0.9}{%
	\makecell[{{p{\linewidth}}}]{% 
	  \tt {\tiny [A$\rangle$GM]}  
	 RAGIONAMENTO STRATEGICO: {'Verifico il costo della loro nuova proposta: A71 12, C17 89, B15 115, A94 118, B96 319, B69 326, A65 351, C82 502, C72 504, A46 514, A28 526, C10 562, A39 568, C79 598, B64 619, A88 621, B19 655, B16 662, C83 830, B23 863, C51 771. Somma progressiva: 12+89=101; +115=216; +118=334; +319=653; +326=979; +351=1330; +502=1832; +504=2336; +514=2850; +526=3376; +562=3938; +568=4506; +598=5104; +619=5723; +621=6344; +655=6999; +662=7661; +830=8491; +863=9354; +771=10125. Supera il limite 9322 di 803: è invalida. Propongo una controproposta che conservi i loro due inserimenti chiave (B23, C51) e C83, ma rientri nel limite massimizzando il mio valore. Per far spazio, rimuovo tre elementi medi con buon rapporto ma sacrificabili per me: A94(118), B69(326), A65(351) totalizzano 795 di sforzo; serve liberarene 803, quindi aggiungo A71(12). Rimuovendo questi 4 si libera esattamente 807, riportando sotto 9322. Parto dalla mia base precedente con C83 incluso (sforzo 8931) e aggiungo B23(863) e C51(771)=+1634 arrivando a 10565; tolgo A94(118), B69(326), A65(351), A71(12)=807, nuovo totale 9758? Attenzione: 10565-807=9758, ancora sopra. Devo togliere ulteriormente 436. Tolgo B15(115) e C17(89) e A25 non è in lista; invece posso togliere C82(502) che per me conta ma è una delle meno essenziali tra i medi: 9758-502=9256; ancora  -? limite 9322, 9256 è sotto di 66, ok. Ora elenco finale e controllo valore per me: accetto la perdita pur di tenere B23, C51, C83 come richiesto. Provo questa proposta equilibrata.'}\\ \tt \\ \tt RIFIUTO: {'A71', 'C17', 'B15', 'A94', 'B96', 'B69', 'A65', 'C82', 'C72', 'A46', 'A28', 'C10', 'A39', 'C79', 'B64', 'A88', 'B19', 'B16', 'C83', 'B23', 'C51'}\\ \tt \\ \tt PROPOSTA: {'B15', 'A94', 'B96', 'A65', 'C72', 'A46', 'A28', 'C10', 'A39', 'C79', 'B64', 'A88', 'B19', 'B16', 'C83', 'B23', 'C51'}\\ \tt \\ \tt ARGOMENTAZIONE: {'La tua proposta supera il limite: ho costruito una versione che mantiene i tre inserimenti chiave per te (B23, C51, C83) insieme all’ossatura condivisa, restando entro 9322. Ho rimosso alcune voci medie per creare spazio senza sacrificare eccessivamente l’impatto complessivo. Se per te è preferibile riaggiungere C82 o B69, posso valutare di sostituirli a coppie con uno tra B15/A94/A65, mantenendo il vincolo. Dimmi quali due elementi tra (C82, B69, A71, C84) valorizzi di più e posso ricalibrare di conseguenza.'} 
	  } 
	   } 
	   } 
	 & & \\ 
 

    \theutterance \stepcounter{utterance}  

    & & \multicolumn{4}{p{0.6\linewidth}}{\cellcolor[rgb]{0.9,0.9,0.9}{%
	\makecell[{{p{\linewidth}}}]{% 
	  \tt {\tiny [GM$\rangle$B]}  
	 RIFIUTO: {'A71', 'C17', 'B15', 'A94', 'B96', 'B69', 'A65', 'C82', 'C72', 'A46', 'A28', 'C10', 'A39', 'C79', 'B64', 'A88', 'B19', 'B16', 'C83', 'B23', 'C51'}\\ \tt \\ \tt PROPOSTA: {'B15', 'A94', 'B96', 'A65', 'C72', 'A46', 'A28', 'C10', 'A39', 'C79', 'B64', 'A88', 'B19', 'B16', 'C83', 'B23', 'C51'}\\ \tt \\ \tt ARGOMENTAZIONE: {'La tua proposta supera il limite: ho costruito una versione che mantiene i tre inserimenti chiave per te (B23, C51, C83) insieme all’ossatura condivisa, restando entro 9322. Ho rimosso alcune voci medie per creare spazio senza sacrificare eccessivamente l’impatto complessivo. Se per te è preferibile riaggiungere C82 o B69, posso valutare di sostituirli a coppie con uno tra B15/A94/A65, mantenendo il vincolo. Dimmi quali due elementi tra (C82, B69, A71, C84) valorizzi di più e posso ricalibrare di conseguenza.'} 
	  } 
	   } 
	   } 
	 & \\ 
 

    \theutterance \stepcounter{utterance}  

    & & & \multicolumn{4}{p{0.6\linewidth}}{\cellcolor[rgb]{1,0.85,0.72}{%
	\makecell[{{p{\linewidth}}}]{% 
	  \tt {\tiny [GM$\langle$B]}  
	 RAGIONAMENTO STRATEGICO: {'La loro ultima proposta include i miei tre ottimi (B23, C51, C83). Voglio spingere per aggiungere almeno uno tra B81/B37/A96/A75/C60/A87, ma rischio un nuovo rifiuto. Prima verifico che la loro proposta sia entro 9322 e valuto piccoli scambi. Calcolo rapido dello sforzo della loro proposta: B15(115)+A94(118)+B96(319)+A65(351)+C72(504)+A46(514)+A28(526)+C10(562)+A39(568)+C79(598)+B64(619)+A88(621)+B19(655)+B16(662)+C83(830)+B23(863)+C51(771) = 9,796? No, ricontrollo: sommo accuratamente. 115+118=233; +319=552; +351=903; +504=1407; +514=1921; +526=2447; +562=3009; +568=3577; +598=4175; +619=4794; +621=5415; +655=6070; +662=6732; +830=7562; +863=8425; +771=9196. Totale 9196, sotto 9322. Buono. Posso accettare oppure chiedere un'aggiunta piccola (ad es. A71=12 o C17=89) per usare capacità residua 126. Inserire A71 e C17 aumenterebbe sforzo di 101, portando a 9297 ≤ 9322: ottimo valore per me, soprattutto C17 non è forte ma il costo è basso; A71 pure. Proporrò l'aggiunta di A71 e C17, lasciando aperto ad accettare la loro proposta se non vogliono. Se rifiutano, posso comunque accettare la loro per chiudere con alto valore.'}\\ \tt \\ \tt ARGOMENTAZIONE: {'La tua proposta è entro il limite e va nella direzione giusta: grazie per aver incluso B23, C51 e C83. Abbiamo ancora circa 126 punti di margine di sforzo: propongo di sfruttarlo aggiungendo due elementi leggeri che non alterano l’equilibrio e aumentano l’impatto complessivo.'}\\ \tt \\ \tt PROPOSTA: {'B15', 'A94', 'B96', 'A65', 'C72', 'A46', 'A28', 'C10', 'A39', 'C79', 'B64', 'A88', 'B19', 'B16', 'C83', 'B23', 'C51', 'A71', 'C17'} 
	  } 
	   } 
	   } 
	  \\ 
 

    \theutterance \stepcounter{utterance}  

    & & \multicolumn{4}{p{0.6\linewidth}}{\cellcolor[rgb]{0.9,0.9,0.9}{%
	\makecell[{{p{\linewidth}}}]{% 
	  \tt {\tiny [A$\langle$GM]}  
	 ARGOMENTAZIONE: {'La tua proposta è entro il limite e va nella direzione giusta: grazie per aver incluso B23, C51 e C83. Abbiamo ancora circa 126 punti di margine di sforzo: propongo di sfruttarlo aggiungendo due elementi leggeri che non alterano l’equilibrio e aumentano l’impatto complessivo.'}\\ \tt \\ \tt PROPOSTA: {'B15', 'A94', 'B96', 'A65', 'C72', 'A46', 'A28', 'C10', 'A39', 'C79', 'B64', 'A88', 'B19', 'B16', 'C83', 'B23', 'C51', 'A71', 'C17'} 
	  } 
	   } 
	   } 
	 & \\ 
 

    \theutterance \stepcounter{utterance}  

    & \multicolumn{4}{p{0.6\linewidth}}{\cellcolor[rgb]{0.8,1,0.9}{%
	\makecell[{{p{\linewidth}}}]{% 
	  \tt {\tiny [A$\rangle$GM]}  
	 RAGIONAMENTO STRATEGICO: {'Verifico lo sforzo della loro proposta aggiungendo A71(12) e C17(89) alla mia ultima lista. Prima lista aveva sforzo: calcolo. Oggetti: B15(115), A94(118), B96(319), A65(351), C72(504), A46(514), A28(526), C10(562), A39(568), C79(598), B64(619), A88(621), B19(655), B16(662), C83(830), B23(863), C51(771). Somma: 115+118=233; +319=552; +351=903; +504=1407; +514=1921; +526=2447; +562=3009; +568=3577; +598=4175; +619=4794; +621=5415; +655=6070; +662=6732; +830=7562; +863=8425; +771=9196. Quindi avevamo 9196, margine 126. Aggiungono A71(12) e C17(89) = +101, nuovo totale 9297, che è entro 9322 con margine 25. Valore per me aggiunto: A71(108) e C17(185) = +293, ottimo. Questa proposta è favorevole e valida. Accetto subito per chiudere con alto punteggio.'}\\ \tt \\ \tt ACCORDO: {'B15', 'A94', 'B96', 'A65', 'C72', 'A46', 'A28', 'C10', 'A39', 'C79', 'B64', 'A88', 'B19', 'B16', 'C83', 'B23', 'C51', 'A71', 'C17'}\\ \tt \\ \tt ARGOMENTAZIONE: {'Accetto: la tua proposta resta entro il limite (sforzo totale 9297) e aumenta l’impatto complessivo con due aggiunte leggere, mantenendo i tuoi elementi chiave. Ottimo compromesso e chiusura efficiente.'} 
	  } 
	   } 
	   } 
	 & & \\ 
 

    \theutterance \stepcounter{utterance}  

    & & & \multicolumn{2}{p{0.3\linewidth}}{\cellcolor[rgb]{0.95,0.95,0.95}{%
	\makecell[{{p{\linewidth}}}]{% 
	  \tt {\tiny [GM$|$GM]}  
	 game successful 
	  } 
	   } 
	   } 
	 & & \\ 
 

    \theutterance \stepcounter{utterance}  

    & & & \multicolumn{2}{p{0.3\linewidth}}{\cellcolor[rgb]{0.95,0.95,0.95}{%
	\makecell[{{p{\linewidth}}}]{% 
	  \tt {\tiny [GM$|$GM]}  
	 end game 
	  } 
	   } 
	   } 
	 & & \\ 
 

\end{supertabular}
}

\end{document}
